% compile: xelatex AUN3_7015101.tex
\documentclass{mko3}
\begin{document}
\mkothreetitle

% ----- หมวดที่ 1 ข้อมูลทั่วไป (รูปแบบ มคอ.3 AUN-QA) -----
\coursecode{7015101}
\coursename{เทคโนโลยีอุบัติใหม่ทางวิศวกรรมคอมพิวเตอร์และปัญญาประดิษฐ์}
\program{หลักสูตรวิศวกรรมศาสตรมหาบัณฑิต สาขาวิชาวิศวกรรมคอมพิวเตอร์และปัญญาประดิษฐ์}
\coursecredits{3 (2-2-5)}
\coursegroup{วิชาเฉพาะด้านบังคับ}
\yearlevel{นักศึกษาปริญญาโท — ภาคปลาย (ตามแผนการศึกษา)}
\semester{2 / 2568}
\instructor{ผศ.ดร.พิศิษฐ์ นาคใจ และ ผศ.ดร.พัชรี มณีรัตน์}
\contact{pisit.nakjai@uru.ac.th, patcharee.manee@uru.ac.th}
\coursedesc{การศึกษาด้านคอมพิวเตอร์หรือปัญญาประดิษฐ์ที่น่าสนใจ และสอดคลอ้งกับเทคโนโลยีปัจจุบัน Study of interesting computer science or artificial intelligence topics aligned with current technologies.}
\courseinfotable
\begin{clodomaintable}
  \cloitem{CLO1}{อธิบาย แนวคิดและหลักการของเทคโนโลยีอุบัติใหม่ทางวิศวกรรมคอมพิวเตอร์และปัญญาประดิษฐ์ ได้อย่างถูกต้อง}{Cognitive}{2}
  \cloitem{CLO2}{ประยุกต์ใช้ เทคโนโลยีอุบัติใหม่ ในการแก้ปัญหาทางวิศวกรรมคอมพิวเตอร์และปัญญาประดิษฐ์}{Cognitive}{3}
  \cloitem{CLO3}{รับรู้ ผลกระทบและการเปลี่ยนแปลงของเทคโนโลยีอุบัติใหม่ ที่มีต่อสังคมและวิชาชีพ}{Affective}{1}
  \cloitem{CLO4}{ตอบสนอง ต่อการเรียนรู้เทคโนโลยีใหม่ ด้วยความกระตือรือร้นและมีส่วนร่วม}{Affective}{2}
  \cloitem{CLO5}{ปฏิบัติตาม แนวทางการใช้เทคโนโลยีอุบัติใหม่ ภายใต้การแนะนำได้อย่างเหมาะสม}{Psychomotor}{3}
\end{clodomaintable}

\begin{cloplotable}{5}{ & \textbf{PLO1} & \textbf{PLO2} & \textbf{PLO3} & \textbf{PLO4} & \textbf{PLO5}}
  \textbf{CLO1} & X &   &   &   &   \\ \hline
  \textbf{CLO2} & X &   &   &   &   \\ \hline
  \textbf{CLO3} &   & X &   &   &   \\ \hline
  \textbf{CLO4} & X &   &   &   &   \\ \hline
  \textbf{CLO5} & X &   &   &   &   \\ \hline
\end{cloplotable}

\begin{bloomtable}
  \bloomrow{CLO1}{2}{}{}{}
  \bloomrow{CLO2}{3}{}{}{}
  \bloomrow{CLO3}{}{1}{}{}
  \bloomrow{CLO4}{}{2}{}{}
  \bloomrow{CLO5}{}{}{3}{3}
\end{bloomtable}

\begin{contentanalysistable}
  \contentrow{การศึกษาด้านคอมพิวเตอร์หรือปัญญาประดิษฐ์ที่น่าสนใจ และสอดคล้องกับเทคโนโลยีปัจจุบัน}{Cognitive 2}{อธิบาย}{ได้อย่างถูกต้อง (U) เพื่อให้เข้าใจหลักการพื้นฐานของเทคโนโลยีอุบัติใหม่ (CLO1) [PLO1]}{4}
  \contentrow{เทคโนโลยีปัจจุบัน}{Cognitive 3}{ประยุกต์ใช้}{ในการแก้ปัญหาทางวิศวกรรมคอมพิวเตอร์และปัญญาประดิษฐ์ได้อย่างเหมาะสม (Ap) เพื่อสร้างนวัตกรรมใหม่ (CLO2) [PLO1]}{4}
  \contentrow{เทคโนโลยีอุบัติใหม่}{Affective 1}{รับรู้}{ผลกระทบที่มีต่อสังคมและวิชาชีพ (Re) เพื่อเตรียมพร้อมรับการเปลี่ยนแปลง (CLO3) [PLO2]}{4}
  \contentrow{การเรียนรู้เทคโนโลยีใหม่}{Affective 2}{ตอบสนอง}{ด้วยความกระตือรือร้นและมีส่วนร่วม (Rs) ในการพัฒนาตนเองอย่างต่อเนื่อง (CLO4) [PLO1]}{4}
  \contentrow{แนวทางการใช้เทคโนโลยีอุบัติใหม่}{Psychomotor 3}{ปฏิบัติตาม}{ภายใต้การแนะนำได้อย่างเหมาะสม (P) เพื่อการนำไปใช้งานจริง (CLO5) [PLO1]}{4}
\end{contentanalysistable}

\mkocoursedesc

\begin{courseactivities}
  \actitem{ศึกษาด้วยตนเอง วิเคราะห์กรณีศึกษา และนำเสนอความเข้าใจ (CLO1)}
  \actitem{ฝึกปฏิบัติ พัฒนาโครงงาน และแก้ปัญหาจริง (CLO2)}
  \actitem{อภิปรายกลุ่ม แสดงความคิดเห็น และเขียนบทสะท้อนคิด (CLO3, CLO4)}
  \actitem{ฝึกปฏิบัติตามคำแนะนำ และทดลองใช้เทคโนโลยี (CLO5)}
\end{courseactivities}

\begin{teachingtable}
  \teachrow{CLO1}{บรรยายแบบมีปฏิสัมพันธ์ - นำเสนอกรณีศึกษา - อภิปรายกลุ่ม [IEL: ประสบการณ์เชิงบูรณาการ ลงมือปฏิบัติ พัฒนาท้องถิ่น]}{การสอบข้อเขียน - การนำเสนอ - รายงานการวิเคราะห์}
  \teachrow{CLO2}{สาธิตการใช้เทคโนโลยี - Project-based Learning - ให้คำปรึกษาโครงงาน [IEL: ประสบการณ์เชิงบูรณาการ ลงมือปฏิบัติ พัฒนาท้องถิ่น]}{ผลงานโครงงาน - การสอบปฏิบัติ - การนำเสนอผลงาน}
  \teachrow{CLO3}{เชิญวิทยากรบรรยาย - จัดกิจกรรมแลกเปลี่ยนความคิดเห็น - อภิปรายผลกระทบต่อสังคม [IEL: ประสบการณ์เชิงบูรณาการ ลงมือปฏิบัติ พัฒนาท้องถิ่น]}{การมีส่วนร่วมในชั้นเรียน - บทความสะท้อนคิด}
  \teachrow{CLO4}{เชิญวิทยากรบรรยาย - จัดกิจกรรมแลกเปลี่ยนความคิดเห็น - อภิปรายผลกระทบต่อสังคม [IEL: ประสบการณ์เชิงบูรณาการ ลงมือปฏิบัติ พัฒนาท้องถิ่น]}{การมีส่วนร่วมในชั้นเรียน - บทความสะท้อนคิด}
  \teachrow{CLO5}{สาธิตขั้นตอนการใช้งาน - ให้คำแนะนำระหว่างปฏิบัติ [IEL: ประสบการณ์เชิงบูรณาการ ลงมือปฏิบัติ พัฒนาท้องถิ่น]}{การทดสอบปฏิบัติ - การสังเกตพฤติกรรม}
\end{teachingtable}

\begin{assessmenttable}
  \assessrow{สอบก่อนเรียน (ถ้ามี)}{---}{0}
  \assessrow{สอบระหว่างเรียน}{CLO1, CLO2, CLO3}{30}
  \assessrow{สอบปลายภาค}{CLO1, CLO2, CLO3, CLO4, CLO5}{40}
  \assessrow{รายงาน / โครงงาน / กิจกรรม}{CLO3, CLO4, CLO5}{30}
\end{assessmenttable}

\begin{weeklyplan}
  \weekrow{1}{บทนำเทคโนโลยีอุบัติใหม่}{บรรยาย + ปฐมนิเทศ}{CLO1}{สอบก่อนเรียน}
  \weekrow{2}{AI และ Machine Learning ยุคใหม่}{บรรยาย + อภิปราย}{CLO1, CLO2}{แบบฝึกหัด}
  \weekrow{3}{IoT และ Edge Computing}{บรรยาย + Workshop}{CLO2}{แบบฝึกหัด}
  \weekrow{4}{Blockchain และ Web3}{บรรยาย + สาธิต}{CLO1}{รายงานย่อย}
  \weekrow{5}{Quantum และเทคโนโลยีอนาคต}{Workshop + อภิปราย}{CLO1}{สอบกลางภาค 1}
  \weekrow{6}{Cloud-native และ DevOps}{บรรยาย + Lab}{CLO2}{แบบฝึกหัด}
  \weekrow{7}{Cybersecurity ยุคใหม่}{บรรยาย + กรณีศึกษา}{CLO3}{แบบฝึกหัด}
  \weekrow{8}{ผลกระทบต่อสังคมและจริยธรรม}{อภิปราย + สะท้อนคิด}{CLO3}{รายงานกลางภาค}
  \weekrow{9}{การประยุกต์ในภาคอุตสาหกรรม}{เชิญวิทยากร + อภิปราย}{CLO2, CLO4}{สอบกลางภาค 2}
  \weekrow{10}{การประยุกต์ในชุมชนท้องถิ่น}{Workshop โจทย์ท้องถิ่น (IEL)}{CLO2, CLO5}{---}
  \weekrow{11}{ออกแบบโครงงาน}{PBL + ให้คำปรึกษา}{CLO5}{ข้อเสนอโครงงาน}
  \weekrow{12}{พัฒนาโครงงาน (ช่วงที่ 1)}{Workshop}{CLO5}{รายงานความก้าวหน้า}
  \weekrow{13}{พัฒนาโครงงาน (ช่วงที่ 2)}{Workshop}{CLO4, CLO5}{---}
  \weekrow{14}{ทดสอบและปรับปรุง}{สาธิต + ทดสอบ}{CLO5}{สาธิตระบบ}
  \weekrow{15}{นำเสนอโครงงานและสรุปรายวิชา}{นำเสนอ + สอบปลายภาค}{CLO1--CLO5}{สอบปลายภาค + โครงงาน}
\end{weeklyplan}

% ----- หมวดที่ 5 Rubric -----
\begin{rubricsection}
\begin{subrubric}{CLO1}{อธิบาย แนวคิดและหลักการของเทคโนโลยีอุบัติใหม่ทางวิศวกรรมคอมพิวเตอร์และปัญญาประดิษฐ์ ได้อย่างถูกต้อง}
  \rubricrow{4 (ดีเยี่ยม)}{อธิบายหลักการได้อย่างถูกต้องลึกซึ้ง เชื่อมโยงทฤษฎีกับบริบทวิชาชีพและวรรณกรรมที่เกี่ยวข้องได้อย่างเหมาะสม}
  \rubricrow{3 (ดี)}{อธิบายหลักการได้ถูกต้องส่วนใหญ่ ครอบคลุมประเด็นสำคัญตามมาตรฐานหลักสูตรปริญญาโท}
  \rubricrow{2 (พอใช้)}{อธิบายได้บางส่วน ยังขาดความลึกหรือการเชื่อมโยงเชิงวิชาการ}
  \rubricrow{1 (ต้องปรับปรุง)}{ไม่สามารถอธิบายได้อย่างถูกต้องหรือไม่ครบถ้วน}
\end{subrubric}

\begin{subrubric}{CLO2}{ประยุกต์ใช้ เทคโนโลยีอุบัติใหม่ ในการแก้ปัญหาทางวิศวกรรมคอมพิวเตอร์และปัญญาประดิษฐ์}
  \rubricrow{4 (ดีเยี่ยม)}{ประยุกต์ใช้ได้อย่างถูกต้องและเหมาะสม แก้ปัญหาซับซ้อนได้ด้วยเหตุผลและหลักฐานสนับสนุน}
  \rubricrow{3 (ดี)}{ประยุกต์ใช้ได้ถูกต้องส่วนใหญ่ ตรงตามวัตถุประสงค์ของงาน}
  \rubricrow{2 (พอใช้)}{ประยุกต์ใช้ได้บางส่วน ยังไม่เหมาะสมหรือไม่สมบูรณ์}
  \rubricrow{1 (ต้องปรับปรุง)}{ไม่สามารถประยุกต์ใช้ได้หรือใช้ผิดแนวทาง}
\end{subrubric}

\begin{subrubric}{CLO3}{รับรู้ ผลกระทบและการเปลี่ยนแปลงของเทคโนโลยีอุบัติใหม่ ที่มีต่อสังคมและวิชาชีพ}
  \rubricrow{4 (ดีเยี่ยม)}{รับรู้ผลกระทบและประเด็นสำคัญได้อย่างลึกซึ้ง เชื่อมโยงกับบริบทวิชาชีพ}
  \rubricrow{3 (ดี)}{รับรู้ได้เหมาะสมส่วนใหญ่}
  \rubricrow{2 (พอใช้)}{รับรู้ได้บางส่วน ยังไม่ครอบคลุม}
  \rubricrow{1 (ต้องปรับปรุง)}{ไม่รับรู้หรือรับรู้ไม่ถูกต้อง}
\end{subrubric}

\begin{subrubric}{CLO4}{ตอบสนอง ต่อการเรียนรู้เทคโนโลยีใหม่ ด้วยความกระตือรือร้นและมีส่วนร่วม}
  \rubricrow{4 (ดีเยี่ยม)}{ตอบสนองต่อการเรียนรู้ใหม่ได้อย่างกระตือรือร้น มีส่วนร่วมและขยายความรู้อย่างต่อเนื่อง}
  \rubricrow{3 (ดี)}{ตอบสนองได้ดี มีส่วนร่วมสม่ำเสมอ}
  \rubricrow{2 (พอใช้)}{ตอบสนองได้บางครั้ง}
  \rubricrow{1 (ต้องปรับปรุง)}{ไม่ตอบสนองหรือมีส่วนร่วมน้อย}
\end{subrubric}

\begin{subrubric}{CLO5}{ปฏิบัติตาม แนวทางการใช้เทคโนโลยีอุบัติใหม่ ภายใต้การแนะนำได้อย่างเหมาะสม}
  \rubricrow{4 (ดีเยี่ยม)}{ปฏิบัติตามแนวทางได้อย่างเหมาะสม มีความรับผิดชอบและปรับปรุงงานได้}
  \rubricrow{3 (ดี)}{ปฏิบัติตามแนวทางได้ดีส่วนใหญ่}
  \rubricrow{2 (พอใช้)}{ปฏิบัติได้บางส่วน ยังต้องกำกับ}
  \rubricrow{1 (ต้องปรับปรุง)}{ไม่สามารถปฏิบัติตามแนวทางได้}
\end{subrubric}
\end{rubricsection}


\begin{learningmaterials}
  \matitem{เอกสารประกอบการบรรยายและบทความล่าสุด}
  \matitem{ตำราและรายงานเทรนด์เทคโนโลยี (Gartner, IEEE Spectrum)}
  \matitem{เครื่องมือทดลอง: Python, cloud platforms, IoT kits}
  \matitem{กรณีศึกษาจากภาคอุตสาหกรรมและชุมชนท้องถิ่น}
  \matitem{สื่อออนไลน์ผ่าน LMS}
\end{learningmaterials}

\begin{supporttable}
  \supportrow{ห้องเรียน Smart Classroom}{1 ห้อง}{พร้อมใช้}{ทันสมัย}
  \supportrow{ห้องปฏิบัติการคอมพิวเตอร์}{1 ห้อง}{พร้อมใช้}{ทันสมัย}
  \supportrow{GPU Server / Cloud Computing Resources}{ครบถ้วน}{พร้อมใช้}{ทันสมัย}
  \supportrow{เอกสารประกอบการสอน, ตำรา, วารสารวิชาการ}{ครบถ้วน}{พร้อมใช้}{ทันสมัย}
\end{supporttable}

\signatureblock{ผศ.ดร.พิศิษฐ์ นาคใจ และ ผศ.ดร.พัชรี มณีรัตน์}{ผศ.ดร.พิทักษ์ คล้ายชม}

\end{document}